\documentclass{classes/progress_report}

% 未指定の項目は原本どおり空欄になる
\author{氏名}
\studentid{学籍番号}
\laboratory{所属研究室}
\reportdate{2026年10月8日}
\reportweek{1}
\reportmode{■対面\hspace{0.5em}□オンタイム遠隔\hspace{0.5em}□オンデマンド}
\researchtheme{研究テーマ}
\advisor{指導担当者}

% 提出前チェック (\prechecked{1,3}で1,3番目を■にする。未指定ならすべて□)
\prechecked{1,2,3,4}
\reviewer{確認者名}
\reviewdate{2026年10月8日}

\begin{document}
\maketitle

\begin{reportbox}{1．前回からの進捗・前回の指示事項への対応}{23mm}
    \hint{前回決めたことに対して、どこまで進んだか。未完了の場合は理由も記入する。}
\end{reportbox}

\begin{reportbox}{2．今回の目標}{17mm}
    \hint{今回の終了時点で何を明らかにする／何を完了する予定だったか。}
\end{reportbox}

\begin{reportbox}{3．今回実施したこと}{38mm}
    \hint{調査・実験・解析・実装・環境構築等を、対象、方法、条件、使用資料・データ・コードが分かるように記入する。}
\end{reportbox}

\begin{reportbox}{4．得られた結果・気づき}{38mm}
    \hint{得られた結果、確認できたこと、想定との差、解釈を記入する。数値・図表・ファイル名等を示すとよい。}
\end{reportbox}

\begin{reportbox}{5．未解決の課題・質問}{38mm}
    \hint{うまくいかなかった点、原因候補、試した対処、教員・上級生へ確認したいことを記入する。}
\end{reportbox}

\begin{reportbox}{6．次回までの課題}{40mm}
    \hint{具体的な作業、完了条件、期限を記入する。複数ある場合は優先順位を付ける。}
\end{reportbox}

\begin{reportbox}{7．上級生・教員からの指示事項／フィードバック}{38mm}
    \hint{学生本人が、助言内容と次に反映する点を記録する。}
\end{reportbox}

\makefooter
\end{document}
